\begin{codebox}
\codeline{\textcolor{cmtgray}{\#\ ==============================}}
\codeline{\textcolor{cmtgray}{\#\ Step\ 1.\ 确定领域与范围}}
\codeline{\textcolor{cmtgray}{\#\ ==============================}}
\codeline{\textcolor{cmtgray}{\#\ 用四个问题界定：}}
\codeline{\textcolor{cmtgray}{\#\ \ \ 本体覆盖什么领域？\ \ \ \ \ \ \ \ ——\ 离散制造车间的设备与加工能力}}
\codeline{\textcolor{cmtgray}{\#\ \ \ 本体用来做什么？\ \ \ \ \ \ \ \ \ \ ——\ 设备调度、能力匹配、冲突检测}}
\codeline{\textcolor{cmtgray}{\#\ \ \ 本体要回答哪些问题？\ \ \ \ \ \ ——\ 能力问题清单（见\ competency-questions.txt）}}
\codeline{\textcolor{cmtgray}{\#\ \ \ 谁使用和维护本体？\ \ \ \ \ \ \ \ ——\ 工艺工程师建模，调度系统消费}}
\codeblank{}
\codeline{\textcolor{cmtgray}{\#\ 范围控制原则：CQ覆盖不到的概念不建模}}
\codeline{\textcolor{cmtgray}{\#\ （例：本体不需要覆盖财务折旧，除非有对应CQ）}}
\codeblank{}
\end{codebox}
\color{black}
